\section{Modelo Dinámico}
Para representar el comportamiento del sistema ante eventos externos, se definen los diagramas de secuencia para los casos de uso críticos.

\subsection{Diagrama de Secuencia: Registro y Clasificación}
Este modelo describe la interacción entre el Frontend (Angular), el API Gateway (Nginx), el Microservicio de Gestión (Spring Boot) y el Microservicio de IA (Python).

\subsection{Vistas de cada perfil}

\subsubsection{Perfil quejoso}
Este perfil está relacionado con la comunidad politécnica; toda aquella persona perteneciente al instituto, al presentar una queja, se convierte directamente en un 'quejoso'.

% --- MQ-01 ---
\begin{figure}[h]
\centering
\includegraphics[width=0.65\textwidth]{imagenes/quejoso/MQ-01_PantallaInicio.png}
\caption{MQ-01: Pantalla de inicio.}
\label{fig:mq01_inicio}
\end{figure}

\noindent \textbf{Descripción de la interfaz (MQ-01):} \
Busca mantener el estilo inicial que el politécnico presenta en sus páginas, apegándose al diseño de la defensoría con un menú de navegación informativo y accesos directos a las acciones principales como Presentar una queja, Seguimiento de queja e Inicar sesión.
\clearpage
% --- MQ-02 ---
\begin{figure}[h]
\centering
\includegraphics[width=0.65\textwidth]{imagenes/quejoso/MQ-02_FormularioRegistro.png}
\caption{MQ-02: Formulario para el registro de quejas.}
\label{fig:mq02_formulario}
\end{figure}

\noindent \textbf{Descripción de la interfaz (MQ-02):} \
Formulario para presentar una queja, abarcando datos personales (nombre, primer y segundo apellido, correo y fecha de nacimiento), información relacionada a la pertenencia al instituto (Unidad académica o Unidad administrativa, Número de boleta), relato de hechos, fecha de hechos, nombre del acusado y carga de evidencias y de identificación.

\clearpage
% --- MQ-03 ---
\begin{figure}[h]
\centering
\includegraphics[width=0.65\textwidth]{imagenes/quejoso/MQ-03_VentanaTutor.png}
\caption{MQ-03: Ventana emergente de información del tutor para menores de edad.}
\label{fig:mq03_tutor}
\end{figure}

\noindent \textbf{Descripción de la interfaz (MQ-03):} \
En caso de que se haya detectado ser menor de edad (por medio de la fecha de nacimiento), esta pantalla permitira ingresar los datos de identificación del tutor del quejoso, como nombre, parentesco y medios de identificación, así como adjuntar una identificación.

\clearpage
% --- MQ-04 ---
\begin{figure}[h]
\centering
\includegraphics[width=0.85\textwidth]{imagenes/quejoso/MQ-04_ExitoFolio.png}
\caption{MQ-04: Pantalla de éxito de generación de folio.}
\label{fig:mq04_exito}
\end{figure}

\noindent \textbf{Descripción de la interfaz (MQ-04):} \\
Pantalla de confirmación final que presenta visualmente el número de folio asignado al trámite. Contiene un botón principal para descargar el acuse de recibo y una sección informativa que invita al usuario a formalizar una cuenta institucional. Ofrece dos propuestas de salida: un botón de acción para iniciar el registro y un enlace de texto para continuar visualizando el trámite como invitado.

\clearpage

% --- MQ-05 ---
\begin{figure}[h]
\centering
\includegraphics[width=0.85\textwidth]{imagenes/quejoso/MQ-05_ConfiguracionCuenta.png}
\caption{MQ-05: Configuración de cuenta institucional.}
\label{fig:mq05_config_cuenta}
\end{figure}

\noindent \textbf{Descripción de la interfaz (MQ-05):} \\
Interfaz de registro que muestra un indicador de progreso de tres pasos (Paso 3: Activar Cuenta). Presenta campos de entrada de texto para el correo electrónico y la creación de una contraseña, incluyendo una guía visual sobre los requisitos de seguridad (caracteres, mayúsculas y números). El botón principal ejecuta la activación, mientras que un enlace en la parte inferior permite abortar el proceso y regresar al inicio.

\clearpage

% --- MQ-06 ---
\begin{figure}[h]
\centering
\includegraphics[width=0.85\textwidth]{imagenes/quejoso/MQ-06_SeguimientoSinCuenta.png}
\caption{MQ-06: Seguimiento de queja sin cuenta.}
\label{fig:mq06_seguimiento_sin_cuenta}
\end{figure}

\noindent \textbf{Descripción de la interfaz (MQ-06):} \\
Ventana emergente que despliega el <<Detalle de la Queja>> de forma estructurada. A la izquierda se muestra un resumen de la denuncia (asunto, fechas y relato original) y una lista de archivos adjuntos con botones de descarga. A la derecha, se visualiza una línea de tiempo vertical que indica el progreso actual del trámite. Al final de la ventana, destaca un botón para descargar el acuse actualizado con el estatus vigente.

\clearpage

% --- MQ-07 ---
\begin{figure}[h]
\centering
\includegraphics[width=0.85\textwidth]{imagenes/quejoso/MQ-07_InicioSesionQuejoso.png}
\caption{MQ-07: Inicio de sesión quejoso.}
\label{fig:mq07_login_quejoso}
\end{figure}

\noindent \textbf{Descripción de la interfaz (MQ-07):} \\
Portal de acceso que utiliza un formulario para el ingreso de usuario y contraseña. La interfaz incluye elementos de soporte como la opción <<¿Olvidaste tu contraseña?>> y un acceso directo para nuevos registros. En la base del formulario, se presenta un cuadro informativo de ayuda técnica para contactar a la Defensoría en caso de problemas de acceso.

\clearpage

% --- MQ-08 ---
\begin{figure}[h]
\centering
\includegraphics[width=0.85\textwidth]{imagenes/quejoso/MQ-08_OpcionesCuenta.png}
\caption{MQ-08: Opciones para crear cuenta.}
\label{fig:mq08_opciones_cuenta}
\end{figure}

\noindent \textbf{Descripción de la interfaz (MQ-08):} \\
Interfaz de bienvenida al portal de seguimiento que presenta dos tarjetas de opción clara. La primera tarjeta está dirigida a usuarios nuevos (Iniciar nueva queja), mientras que la segunda permite a usuarios con trámites previos activar su perfil digital (Activar cuenta con folio). Es una pantalla de navegación limpia diseñada para segmentar al usuario según su situación actual.

\clearpage

% --- MQ-09 ---
\begin{figure}[h]
\centering
\includegraphics[width=0.85\textwidth]{imagenes/quejoso/MQ-09_CuentaExistente.png}
\caption{MQ-09: Crear cuenta con queja existente.}
\label{fig:mq09_cuenta_existente}
\end{figure}

\noindent \textbf{Descripción de la interfaz (MQ-09):} \\
Formulario de vinculación que solicita tres datos clave: correo electrónico, número de folio y contraseña. La interfaz mantiene un diseño consistente con el registro estándar, añadiendo el campo de folio para validar la propiedad de la queja. Al final, presenta un botón de acción para finalizar la activación y un enlace para retornar al menú anterior.

\clearpage

% --- MQ-10 ---
\begin{figure}[h]
\centering
\includegraphics[width=0.85\textwidth]{imagenes/quejoso/MQ-10_RestablecerPass.png}
\caption{MQ-10: Restablecer contraseña.}
\label{fig:mq10_reset_pass}
\end{figure}

\noindent \textbf{Descripción de la interfaz (MQ-10):} \\
Pantalla minimalista centrada en un único campo de entrada para el correo electrónico institucional. Incluye un breve texto instructivo sobre el envío del enlace de recuperación.

\clearpage

% --- MQ-11 ---
\begin{figure}[h]
\centering
\includegraphics[width=0.85\textwidth]{imagenes/quejoso/MQ-11_GestionTramites.png}
\caption{MQ-11: Gestión de trámites.}
\label{fig:mq11_gestion}
\end{figure}

\noindent \textbf{Descripción de la interfaz (MQ-11):} \\
Panel de control (Dashboard) que ofrece un resumen numérico de las quejas en curso, finalizadas y pendientes a través de tarjetas de indicadores. Debajo se despliega una tabla con el listado de trámites, mostrando folio, fecha, asunto y estatus. Cada registro cuenta con una columna de <<Acciones>> mediante iconos para editar, eliminar o visualizar el detalle. Incluye una barra de búsqueda específica por folio para navegación rápida.

\clearpage

% --- MQ-12 ---
\begin{figure}[h]
\centering
\includegraphics[width=0.85\textwidth]{imagenes/quejoso/MQ-12_ConfirmarEliminacion.png}
\caption{MQ-12: Confirmar eliminación.}
\label{fig:mq12_borrar}
\end{figure}

\noindent \textbf{Descripción de la interfaz (MQ-12):} \\
Ventana modal de advertencia que aparece sobre el tablero de gestión. Presenta un icono de alerta amarillo y un mensaje de confirmación que indica el número de folio específico a eliminar. Contiene dos botones contrastados: <<Cancelar>> para cerrar la ventana sin cambios y <<Eliminar definitivamente>> advirtiendo que esta es irreversible.

\clearpage

% --- MQ-13 ---
\begin{figure}[h]
\centering
\includegraphics[width=0.85\textwidth]{imagenes/quejoso/MQ-13_EditarQueja.png}
\caption{MQ-13: Editar queja.}
\label{fig:mq13_editar}
\end{figure}

\noindent \textbf{Descripción de la interfaz (MQ-13):} \\
Pantalla de edición que retoma el formulario original de la queja. Muestra un recuadro informativo resaltado indicando el estatus actual y las limitaciones de edición vigentes. Permite modificar campos de texto como el asunto y el relato, así como gestionar la sección de evidencias, donde el usuario puede eliminar archivos previos o adjuntar nuevos documentos mediante un botón dedicado.

\clearpage

% --- MQ-14 ---
\begin{figure}[h]
\centering
\includegraphics[width=0.85\textwidth]{imagenes/quejoso/MQ-14_VerQueja.png}
\caption{MQ-14: Ver queja.}
\label{fig:mq14_ver}
\end{figure}

\noindent \textbf{Descripción de la interfaz (MQ-14):} \\
Interfaz informativa dividida en dos columnas principales sobre fondo gris claro. La columna izquierda contiene el resumen descriptivo de la denuncia y la lista de evidencias con enlaces de descarga. La columna derecha muestra un indicador visual de progreso del trámite (stepper) con marcas de verificación en los pasos completados. En la parte inferior se ubican botones para descargar el acuse oficial y retornar al tablero principal.

\clearpage

% --- MQ-15 ---
\begin{figure}[h]
\centering
\includegraphics[width=0.85\textwidth]{imagenes/quejoso/MQ-15_HistorialQuejas.png}
\caption{MQ-15: Historial de quejas.}
\label{fig:mq15_historial}
\end{figure}

\noindent \textbf{Descripción de la interfaz (MQ-15):} \\
Pantalla de consulta avanzada que presenta un robusto sistema de filtros en la parte superior, permitiendo búsquedas por folio, estatus y rangos de fecha. La tabla inferior organiza el historial cronológicamente y añade columnas de información como la <<Unidad Académica>> involucrada. Incluye controles de paginación en la base para navegar entre múltiples registros de manera eficiente.

\clearpage

% --- MQ-16 ---
\begin{figure}[h]
\centering
\includegraphics[width=0.85\textwidth]{imagenes/quejoso/MQ-16_NuevaDenuncia.png}
\caption{MQ-16: Nueva denuncia desde cuenta.}
\label{fig:mq16_nueva_denuncia}
\end{figure}

\noindent \textbf{Descripción de la interfaz (MQ-16):} \\
Formulario de denuncia que presenta la información del solicitante de forma pre-llenada y bloqueada para mayor agilidad. Se divide en bloques lógicos: Identificación, Detalles del Hecho y Gestión de Archivos. Destaca un área de <<Drag and Drop>> resaltada para el adjunte de evidencias y botones finales que permiten tanto el envío oficial como el guardado de un borrador para terminar después.

\clearpage

% --- MQ-17 ---
\begin{figure}[h]
\centering
\includegraphics[width=0.85\textwidth]{imagenes/quejoso/MQ-17_FolioCuenta.png}
\caption{MQ-17: Generación de folio de queja enviada desde cuenta.}
\label{fig:mq17_folio_cuenta}
\end{figure}

\noindent \textbf{Descripción de la interfaz (MQ-17):} \\
Pantalla de éxito personalizada para el usuario con sesión iniciada. Mantiene la estructura de confirmación de folio, pero integrada dentro del panel lateral de navegación del usuario. Presenta el icono de éxito en verde y el botón de descarga de acuse, eliminando las opciones de creación de cuenta al ser un usuario ya registrado.

\clearpage

% --- MQ-18 ---
\begin{figure}[h]
\centering
\includegraphics[width=0.85\textwidth]{imagenes/quejoso/MQ-18_Notificaciones.png}
\caption{MQ-18: Centro de notificaciones.}
\label{fig:mq18_notificaciones}
\end{figure}

\noindent \textbf{Descripción de la interfaz (MQ-18):} \\
Espacio de mensajería que utiliza un sistema de pestañas para filtrar entre <<Todas>>, <<Pendientes>> y <<Leídas>>. Los avisos urgentes o requerimientos de información aparecen en recuadros de color rojo con botones de acción inmediata (Completar requerimiento). El resto de notificaciones se muestran en una lista cronológica con indicadores visuales de lectura.

\clearpage

% --- MQ-19 ---
\begin{figure}[h]
\centering
\includegraphics[width=0.85\textwidth]{imagenes/quejoso/MQ-19_PropuestaConciliacion.png}
\caption{MQ-19: Propuesta de conciliación.}
\label{fig:mq19_conciliacion}
\end{figure}

\noindent \textbf{Descripción de la interfaz (MQ-19):} \\
Documento digital interactivo que presenta el acuerdo propuesto por la Defensoría. El usuario visualiza los puntos clave del convenio y dispone de dos botones de gran formato para <<Aceptar>> (verde) o <<Rechazar>> (blanco con borde rojo). En caso de rechazo, se habilita un cuadro de texto inferior para redactar los motivos, junto con una opción para descargar la propuesta en PDF.

\clearpage

% --- MQ-20 ---
\begin{figure}[h]
\centering
\includegraphics[width=0.85\textwidth]{imagenes/quejoso/MQ-20_ConfiguracionPerfil.png}
\caption{MQ-20: Configuración de perfil.}
\label{fig:mq20_perfil}
\end{figure}

\noindent \textbf{Descripción de la interfaz (MQ-20):} \\
Panel de gestión de cuenta organizado en secciones: Identidad Institucional (datos no editables), Información de Contacto y Datos del Tutor. Los campos editables aparecen con bordes definidos para su modificación. Incluye un botón destacado en la esquina superior para el cambio de contraseña y un botón inferior para salvar los cambios realizados en la información de contacto o representantes.

\subsubsection{Perfil recepcionista}
DIAGRAMAS DE  SECUANCI, ACTIVIDADES, CASOS DE USO, MOCKUOS, SE LLAMAN CUANDO SE USEN CAT MENSAJES Y CAT ERRORES


